\documentclass{article}
\usepackage{amsmath,amssymb,amsthm}
\newtheorem{theorem}{Theorem}
\begin{document}
\title{A conditional positive composition for paired Haar readouts}
\maketitle

\paragraph{Status.}
The composition theorem below is a written proof. It does not establish
the common-residue hypothesis for arbitrary child words. The finite
certificates at the end use exact rational arithmetic, not Lean.

For $D>p$ put
\[
 f_t(D,p,a)=\frac{a+t}{D-p}-\frac{a-t}{D+t-p}.
\]
For $0\leq r,s<m$, let the distinct real poles $p_r$ satisfy
$p_r+p_s<1$, and put
\[
 K_{sr}=(1-p_s-p_r)^{-2}.
\]
This matrix is positive definite: it is the Gram matrix of the linearly
independent functions $u\mapsto e^{p_ru}$ in the positive measure
$u e^{-u}\,du$ on $(0,\infty)$. Linear independence follows, for example,
by differentiating at zero and using the Vandermonde determinant.

\begin{theorem}
Suppose two child butterfly words at the same lower-level pole geometry
admit a common output-indexed residue vector $a\in\mathbb R^m$, with
independent node biases and strict signs for each child word.
Then, for any prescribed final-stage switch word
$\theta\in\{0,1\}^m$, the combined parent word is feasible after
introducing a sufficiently small new pole gap $t>0$.
\end{theorem}
\begin{proof}
Write $\sigma_s=1-2\theta_s$, and choose $v=K^{-1}\sigma$.
There is $\delta>0$ so small that replacing the child residues by
$a+\delta v/2$ and $a-\delta v/2$, respectively, preserves every strict
child sign. This follows from continuity and the finiteness of the
child query sets; the original child biases can be retained.

The final switch sends the first child's token $r$ to pole
$p_r+\theta_rt$ and the second child's token $r$ to pole
$p_r+(1-\theta_r)t$. Give these two physical output poles the residues
$a_r+\delta v_r/2$ and $a_r-\delta v_r/2$, respectively. These assignments
are allowed to depend on the parent switch word. At every lower-stage
query, the resulting score converges to its already-correct child score
as $t\downarrow0$.

At the final stage, take the node bias to be zero. Its score at prefix
$s$ is
\[
 T_s(t)=\sum_r\left[
 f_t(1-p_s,p_r+\theta_rt,a_r+\delta v_r/2)
 -f_t(1-p_s,p_r+(1-\theta_r)t,a_r-\delta v_r/2)
 \right].
\]
Using
\[
 \frac{f_t(D,p,a)}{t}
 =\frac{a}{(D-p)(D+t-p)}+
   \frac1{D-p}+\frac1{D+t-p},
\]
one obtains
\[
 \lim_{t\downarrow0}\frac{T_s(t)}t
 =\delta\sum_r K_{sr}v_r=\delta\sigma_s.
\]
Consequently all final-stage signs are correct for sufficiently small
$t>0$. Finitely many positive denominator margins and child sign
margins persist simultaneously. This proves the claim.
\end{proof}

\paragraph{What this does not prove.}
The new gap is selected after the lower geometry and its strict margins
are fixed. No bound such as $t=\varepsilon^{2^d}$ follows. Recursively
much smaller gaps are compatible with the original local full-cube
positivity condition, since that condition only bounds the sum of gaps.
The remaining hypothesis is substantive: it asks for one residue vector
that works for two independently chosen child routing words. Extending
that intersection statement by a common perturbation at the next level
would require simultaneous conjugated-Cauchy positivity systems, and
those can fail. Moreover, a scalar Haar theorem still needs a supported
local-domain construction before it implies a worst-case head bound.

\paragraph{Two exact finite checks of a stronger invariant.}
At $\gamma_k=(1/100)^{2^k}$, one common physical-output-indexed residue
vector realizes all $256$ eight-token words with bottom-stage bits fixed
to zero. Independent node biases are allowed for each word. The file
\texttt{joint-family-haar-common-residue-primal.json} gives a homogeneous
rational vector whose final coordinate is positive; every reduced
strict inequality is positive.

For sixteen tokens, five explicitly listed words have no common residue
vector, even allowing independent node biases for every word. The file
\texttt{joint-family-haar-common-residue-dual.json} contains eighteen
strictly positive rational weights, summing to one, whose weighted
reduced rows sum to zero in all seventeen homogeneous coordinates.
The solver-free verifier
\texttt{JointFamilyHaarCommonResidue.py} checks both certificates.
Each of the ten pairs, ten triples, and five four-element subfamilies of
these five words was separately feasible in exact rational linear
programming. Those exploratory subfamily checks are not needed to
verify the saved certificates. In particular, the five-word obstruction
does not settle the pairwise intersection hypothesis in the theorem.

\end{document}
